\section{Distinct multiscale arguments and finite refinements}
\label{app:dyadic-alternatives}
\subsection{Arbitrary partitions and the dense predecessor}
The exact dyadic formula in Theorem~\ref{thm:dyadic-exact} relies on a
simultaneous prefix-hitting construction. A separate logarithmic estimate
works for arbitrary partitions in \eqref{eq:dyadic-polynomial}.

\begin{proposition}[Partition-independent estimate]
\label{prop:arbitrary-dyadic}
For the polynomial and means in \eqref{eq:dyadic-polynomial}, with any
choice of equal-size level partitions,
\[
 0<H\le5+\log_2(1+(L-1)\log2),\qquad T=L.
\]
\end{proposition}
\begin{proof}
Use the quantile surrogate from \eqref{eq:dyadic-surrogate}, writing
$r(t)$ for the nonincreasing failure-count quantile and $\delta=2^{-L}$.
The interval $0<t\le\delta$ contributes at most
$\delta\sum_{j=1}^L2^j<2$, and $t>1/2$ contributes zero. For
$\delta<t\le1/2$, let $l$ be the largest integer with $t\le2^{-l}$.
Write $\log_2^+v=\max\{0,\log_2v\}$ for $v>0$.
Then $2^l\le1/t$, and for $r>0$,
\begin{equation}\label{eq:dyadic-log-sum}
 \sum_{j=1}^l\min(2^j,r)
 \le r\left(3+\log_2^+\frac1{tr}\right).
\end{equation}
If $r\ge2^l$, the sum is less than $2r$. If $r<2$, it is $lr$ and
$l\le1+\log_2(2^l/r)$. Otherwise, with $h=\lfloor\log_2r\rfloor$,
the first $h$ capped terms sum to less than $2r$ and the others sum to
$(l-h)r$; use $h>\log_2r-1$. This proves
\eqref{eq:dyadic-log-sum}, interpreted as zero at $r=0$.

Put $D=(\delta,1/2)$ and $M_0=\int_Dr(t)\,dt\le1$. As
$\log_2^+v\le\log(1+v)/\log2$, the bound becomes
\[
 H\le2+3M_0+\frac1{\log2}
       \int_Dr(t)\log\left(1+\frac1{tr(t)}\right)\,dt.
\]
For $M_0>0$, Jensen's inequality under the probability measure
$r(t)\,dt/M_0$ gives
\[
 \int_Dr(t)\log\left(1+\frac1{tr(t)}\right)\,dt
 \le M_0\log\left(1+\frac{(L-1)\log2}{M_0}\right)
 \le\log(1+(L-1)\log2).
\]
The final step uses that $M\log(1+B/M)$ increases for $M>0$, since
$\log(1+s)\ge s/(1+s)$. At $M_0=0$ the integral is zero. This proves
the upper bound. Independent rounding and common-threshold rounding
have distinct expected polynomial values at these interior means, so $H>0$.
\end{proof}

A dense symmetric construction gives a different exact reduction.
With $m=2^L$, $u_j=2^{-j}$, $k_j=mu_j$, define
\begin{equation}\label{eq:dense-predecessor}
 \widetilde f_L(a,z)=\sum_{j=1}^L\frac1{u_j\binom m{k_j}}
             \sum_{|K|=k_j}a_j\prod_{i\in K}z_i,
 \qquad a_j=u_j,\quad z_i=1-1/m.
\end{equation}
Again $T=\cav\widetilde f_L=L$. Conditional on $r$ failed leaves, the
fraction of size-$k_j$ subsets hit is
\[
 g_j(r)=1-\frac{\binom{m-r}{k_j}}{\binom m{k_j}}
       \le\min(1,u_jr).
\]
The inequality is the union bound: a uniform size-$k_j$ subset contains
each specified failed leaf with probability $k_j/m=u_j$.
Thus the deficiency is $\sum_j A_jg_j(R)/u_j$, bounded by the same
surrogate $\sum_j A_j\min(2^j,R)$ and consequently by
Proposition~\ref{prop:arbitrary-dyadic}.

In this dense model the hull gap itself is exactly the optimum of the
finite rational program
\begin{equation}\label{eq:dense-count-lp}
 \begin{aligned}
 \max\quad&\sum_{j=1}^L\sum_{r=0}^m z_{jr}g_j(r)/u_j\\
 \text{subject to}\quad& p_r\ge0,\quad 0\le z_{jr}\le p_r,\\
 &\sum_{r=0}^mp_r=1,\quad \sum_{r=0}^mrp_r=1,
       \quad\sum_{r=0}^mz_{jr}=u_j\quad(j=1,\ldots,L).
 \end{aligned}
\end{equation}
Every feasible binary law supplies $p_r=\Prob(R=r)$ and
$z_{jr}=\Prob(A_j=1,R=r)$. Conversely draw $R$ with probabilities $p_r$,
choose a uniform $r$-subset of failed leaves, and, conditional on $R=r$,
choose anchors independently with probabilities $z_{jr}/p_r$ when
$p_r>0$. These choices realize all means and the objective; zero-mass
states can be ignored. Thus \eqref{eq:dense-count-lp} is an exact reduction,
not a formula for the sparse model. No floating-point values from this
program are needed for the divergent bound.

\subsection{Dyadic conditional-failure rounding and its tail lemma}
\label{app:dyadic-rounding}
The following proof uses finitely many dyadic scales rather than a
harmonic kernel. Though its $O(\log d)$ bound is weaker, its tail
argument is distinct.

\begin{lemma}[Tail integral]\label{lem:dyadic-tail}
For a finite collection $p_j\ge0$, let $N(s)=|\{j:p_j\ge s\}|$ and
$S=\sum_jp_j$. For $u>0$,
\begin{equation}\label{eq:dyadic-tail}
 \int_0^u\min\{N(s),u/s\}\,ds\ge\min(u,S).
\end{equation}
\end{lemma}
\begin{proof}
If $S\le u$, then every $p_j\le u$ and $\int_0^uN(s)\,ds=S$.
Monotonicity gives $sN(s)\le\int_0^sN(t)\,dt\le S\le u$, so the
integrand equals $N(s)$ throughout. If $S>u$, then
$\int_0^uN(s)\,ds\ge u$: either a $p_j\ge u$, or that integral equals
$S$. By continuity choose $v\le u$ with $\int_0^vN(s)\,ds=u$.
For $s\le v$ the same monotonicity bound gives $sN(s)\le u$.
Integrating on $(0,v)$ therefore gives $u$.
\end{proof}

\begin{proposition}[Dyadic degree bound]\label{prop:dyadic-degree}
Let $K=1+\lfloor\log_2(d-1)\rfloor$ for $d\ge2$. Every positive
multilinear polynomial of maximum degree at most $d$ on a nonnegative
box satisfies $T\le2(K+1)H/c_0$.
\end{proposition}
\begin{proof}
For each scale $B\in\{1,2,\ldots,2^{K-1}\}$ restricted to powers of
two, draw a common uniform $U$. Low coordinates succeed when $U\le x_i$.
A high coordinate with failure mean $p>0$ has $h=\min(Bp,1)$ and,
conditional on $U$, fails independently with probability $p/h$ when
$U\le h$, and never otherwise. Coordinates with $p=0$ never fail.
The marginal failure is $h(p/h)=p$ and the probabilities are valid
because $B\ge1$.

For a one-low term with anchor mean $u>0$, let $p_1,\ldots,p_r$ be
its high failure means. At $U=t\le u$, each $p_j\ge t/B$ is active
and has conditional failure probability at least $1/B$. Conditional
independence and $1-e^{-z}\ge c_0\min(1,z)$ give deficiency
\[
 G_B\ge c_0\int_0^u\min(1,N(t/B)/B)\,dt
       =c_0\int_0^{u/B}\min(B,N(s))\,ds.
\]
For $0<s\le u$ with $N(s)>0$, choose the largest power of two
$B\le\min(N(s),u/s)$. It is in the allowed set, since
$N(s)\le r\le d-1$, and exceeds half this minimum. Therefore
\[
 \sum_B\ind\{s\le u/B\}\min(B,N(s))
       \ge\tfrac12\min(N(s),u/s).
\]
Integrating and invoking Lemma~\ref{lem:dyadic-tail} yields
$\sum_BG_B\ge c_0\min(u,S)/2=c_0t_e/2$.
Terms with $u=0$ have zero gap. Independence supplies the same guarantee
for terms with zero or at least two low coordinates by
Lemma~\ref{lem:easy-terms}. The uniform mixture of the $K$ scale laws
and independence thus guarantees $c_0t_e/(2(K+1))$ for every term.
Sum and apply Proposition~\ref{prop:box-transfer}.
\end{proof}

\subsection{The simplest harmonic cutoff and its reciprocal correction}
\label{app:reciprocal-cutoff}
The general finite theorem permits $M=d-1$, so $\eta=1$ and
$\Lambda=1+\log(d-1)$. Let $z_\Lambda$ solve
$J_{\Lambda,1}(z_\Lambda)=z_\Lambda$ and put
$\alpha=\Lambda z_\Lambda$, $\ell=\log\Lambda$. For $\ell\ge4$,
\begin{equation}\label{eq:rho-one-finite}
 z_\Lambda\ge\frac{\ell-\log\ell-2}{\Lambda}.
\end{equation}
Indeed $A=\ell-\log\ell-2\ge1/2$ and $A\le\ell<\Lambda$.
Equation~\eqref{eq:lambert-integral-certificate} with $\eta=1$ gives
$\Lambda J(A/\Lambda)\ge\ell-\log\ell-1-1/(2A)\ge A$.
Monotonicity of $J(z)-z$ proves the finite bound.

The integrated linear upper bound and quadratic lower bound imply
\[
 \ell-\log\alpha-1-\frac1{2\alpha}
 \le\alpha\le\ell-\log\alpha-1+\frac\alpha\Lambda.
\]
Together with \eqref{eq:rho-one-finite}, these show $\alpha/\ell\to1$:
once $\alpha\ge1$, the upper bound gives
$\alpha\le\ell/(1-1/\Lambda)$, while the finite lower bound is
$\ell-\log\ell-2$. Substitution yields
\begin{equation}\label{eq:rho-one-asymptotic}
 \Lambda z_\Lambda=\log\Lambda-\log\log\Lambda-1+o(1).
\end{equation}
The constant $-1$ belongs to this fixed cutoff, whereas the tuned
cutoff in \eqref{eq:second-order-upper} removes that constant after
expressing the bound with numerator $N$.

There is also a reciprocal refinement. If $w e^w=\Lambda/e$, then
\begin{equation}\label{eq:rho-one-reciprocal}
 \alpha=w-\frac1{2w}+O(w^{-2})\qquad(\Lambda\to\infty).
\end{equation}
To prove it, integrate the cubic Taylor bound for $1-e^{-x}$ with
$x=(1-v)/(\Lambda v)$. After multiplication by $\Lambda$, the error
beyond the quadratic term is at most
\[
 \Lambda\int_{\alpha/\Lambda}^1\frac{x^3}{6}\,dv
 \le\frac1{12\alpha^2}.
\]
The exact quadratic integral is
$\int_z^1(1/v-1)^2\,dv=1/z+2\log z-z$.
Its omitted terms, and the linear endpoint term, are
$O(\log\Lambda/\Lambda)=o(\alpha^{-2})$ because $\alpha\sim\ell$.
Consequently
\[
 \alpha+\log\alpha=\log\Lambda-1-\frac1{2\alpha}
                       +O(\alpha^{-2}).
\]
Compare this with $w+\log w=\log\Lambda-1$. Since $\alpha\sim w$,
the mean value theorem for $t+\log t$ first gives
$\alpha-w=O(1/w)$ and, after substitution, gives
\eqref{eq:rho-one-reciprocal}. This refines the scalar certificate only;
it is not a second-order lower bound on $R_d$.

\subsection{A second rational two-level parameter choice}
\label{app:two-level-variant}
A related finite variant illustrates that the two-level mechanism is not
tied to the particular coefficients in Theorem~\ref{thm:cubic-two-level}.
For real $a,c\in[0,1]$, direct expansion gives
\begin{equation}\label{eq:two-level-variant-scalar}
\begin{aligned}
 ac^2+\frac{24}{25}a^2
 -\left(\frac{41}{25}a+\frac45c-\frac{68}{75}\right)
 &=\frac{24}{25}\left(a-\frac{41-25c^2}{48}\right)^2\\
 &\quad+\frac{(1-c)(5c-3)^2(5c+11)}{480}.
\end{aligned}
\end{equation}
Both terms on the right are nonnegative. The equality atoms
$(1/3,1)$ and $(2/3,3/5)$, each with probability $1/2$, have mean
$(1/2,4/5)$ and scalar envelope value $83/150$.
For $25\mid m$, use groups of $m$ variables with these two means and
$f_m=A E_2(W)+(24m/25)E_2(U)$. At binary normalized counts,
\[
 \frac{2f_m}{m^3}=ac^2+\frac{24}{25}a^2
                   -\frac{ac+(24/25)a}{m}.
\]
The expected correction is at most $49/(50m)$ because $ac\le a$.
The scaled concave envelope and termwise gap are, respectively,
$(49/50)(1-1/m)$ and $(22/25)(1-1/m)$; cubic terms have lower envelope
$1/10$, while quadratic terms have lower envelope zero. The affine
minorant gives scaled hull gap at most $49/50-83/150=32/75$.
Conditional independent sampling from the equality atoms, as in
Theorem~\ref{thm:cubic-two-level}, gives scaled hull gap at least
$(32/75)(1-1/m)$. Thus
\[
 \frac{33}{16}(1-1/m)\le\frac{T(f_m)}{H(f_m)}\le\frac{33}{16}.
\]
At $m=50$ the lower certificate is $1617/800>2$. This variant has a
smaller limiting ratio than the two-level family in the main text.
